\section{結論}
%----------------------------------------------------------------------

本稿では，連合学習における計算異種性（モデル幅の差異）と
アーキテクチャ異種性（CNN/ViT の混在）を同時に解決する
フレームワーク AFAD を提案した．
32次元共有ボトルネック層を接点として HeteroFL と FedGen を統合する設計により，
幅・アーキテクチャともに多様な端末が単一の FL 環境に参加できる．

OrganAMNIST を用いた IID 環境での実験において，
AFAD はサーバ精度 82.93\% を達成した．
HeteroFL（82.20\%）を上回り，
二重異種性への同時対応が HeteroFL 比での精度低下を伴わないことを実証した．
FedGen（86.32\%）との差 3.4 ポイントは今後の課題として残るが，
HeteroFL が参加を許容しない ViT 端末を含む全 10 台という
より困難な設定でこの結果を達成した点は，提案手法の実用的可能性を示唆する．

今後の課題として，Non-IID 環境での評価，ViT 端末の計算制約のさらなる拡張，
および FitNet スタイルの蒸留による低能力端末の backbone 品質向上を検討する予定である．
AFAD の設計原則が，計算資源の制約を超えて多様な端末が協調学習に参加できる
連合学習の実用化に向けた一助となれば幸いである．

